\section*{Complete author Section 7.1.1 proof}
\paragraph{Proof of Theorem 38.} By using Theorem 37, we obtain that:
\[A(x)=(x^2+1)A(x^2)+\frac{A(x^2)}x.\tag{3}\]
Suppose $B_d(x):=\mathrm{trunc}(A(x),\lfloor d/2\rfloor+1)$. Notice that the degree of $C_d(x):=(x^2+1)B_d(x^2)+B_d(x^2)/x$ is at most $2\lfloor d/2\rfloor+4$ and $\mathrm{trunc}(C_d(x),d)=\mathrm{trunc}(A(x),d)$. Hence we can compute $\mathrm{trunc}(A(x),d)$ from $C_d(x)$ by subtracting at most 4 monomials, which can be done using $O(\log d)$ gates. Also $B_d(x)$ can be computed from $\mathrm{trunc}(A(x),\lfloor d/2\rfloor)$ using $O(\log d)$ gates. Hence we obtain the following recurrence:
\[L(\mathrm{trunc}(A(x),d))\le L(\mathrm{trunc}(A(x),\lfloor d/2\rfloor))+O(\log d).\]
This implies, $L(\mathrm{trunc}(A(x),d))=O(\log^2d)$.\qed
\section*{Complete author Section 7.1.2 proof}
\paragraph{Proof of Theorem 40.} We use the notation $\nu_2(m)$ to denote the highest power of 2 which divides $m\in\mathbb N$. We partition the set $[d]$ into $\lfloor\log d\rfloor$ sets $S_0,S_1,S_2,\ldots,S_{\lfloor\log d\rfloor}$ such that $k\in S_i$ iff $\nu_2(k)=i$. We define the set $O_m:=\{k:k\le m\text{ and }k\text{ is odd}\}$. Now, notice that $S_i=\{2^ik:k\in O_{\lfloor d/2^i\rfloor}\}$. For a set $S\in\mathbb N$, we define the polynomial $g_S:=\sum_{i\in S}x^i$. Observe that:
\[\mathrm{trunc}(F(x),d)=1+\sum_{i=1}^{\lfloor\log d\rfloor}(-1)^ig_{S_i}.\]
Trivially, $g_{S_i}=g_{O_{\lfloor d/2^i\rfloor}}(x^{2^i})$. Also notice that $g_{O_m}=g_{[m]}-g_{\lfloor m/2\rfloor}(x^2)$. Therefore, $L(g_{O_m})=(\log m)$, which implies that $g_{S_i}=O(\log d)$. Hence, $L(\mathrm{trunc}(F(x),d))=O(\log^2d)$.\qed
\paragraph{Remark 41.} Note that, there are power series like $\sum_{i\ge0}x^{i!}$ which are transcendental and their truncations up to degree $d$ are easy to compute. However, the series is highly sparse and degree-$d$ truncations has only $\mathrm{poly}(\log d)$ monomials, hence the easiness is trivial. The examples we discover in this work are of dense power series.
\section*{Complete author Section 7.2 comparison}
Consider the holonomic sequence $f_n=(n!)$ defined by $f_0=1$ and $f_{n+1}-(n+1)f_n=0$. Also consider the corresponding generating function $F(x)=\sum_{n\ge0}n!x^n$. We now show that $F(x)$ is transcendental and that truncation of $F(x)$ is (conditionally) hard to compute.
\paragraph{Corollary 43.} The power series $F(x)=\sum_{n\ge0}n!x^n$ is transcendental.
\begin{proof}It is clear that the radius of convergence of $F(x)$ is zero (follows from the ratio test). Hence Lemma 42 implies that $F(x)$ is transcendental.\end{proof}
\paragraph{Proof of Theorem 44.} We know that $d!x^d=\mathrm{trunc}(F(x),d)-\mathrm{trunc}(F(x),d-1)$. Setting $x=1$, we conclude.\qed
